% docs/texte/05_theory/fundamentals.tex -- theory chapter "Radio Astronomical Fundamentals" (Interferosim)
% Version: 0.1.0 (2026-10-12)
% Notation as defined in texte/05_theory/notation.tex.
% ================================================================

\index{brightness temperature}\index{antenna temperature}
\section{Brightness temperature and flux density}

In the Rayleigh--Jeans limit the specific intensity of a black body is
\begin{equation}
  I_\nu = \frac{2\,k\,T_b}{\lambda^2} .
\end{equation}
The brightness temperature $T_b$ is defined by this relation for any
source, thermal or not. The flux density is the intensity integrated
over the source,
\begin{equation}
  S = \int_{\mathrm{source}} I_\nu\,\mathrm{d}\Omega
    \;\approx\; \frac{2\,k\,T_b\,\Omega}{\lambda^2}
  \qquad\text{(uniform source).}
  \label{eq:flux-tb}
\end{equation}
\emph{Example (notebook~B):} the quiet Sun at 21\,cm has
$T_b \approx 10^5\,\mathrm{K}$ and a radio diameter of about $34'$, which
gives $S \approx 4.8\times10^5\,\mathrm{Jy} \approx 48\,\mathrm{SFU}$.
Both quantities vary with the solar cycle by about a factor of two.

\section{Antennas}

An antenna collects power from a solid angle given by its beam. The
effective area $A_{\mathrm{eff}}$ and the gain $G$ are related by
\begin{equation}
  G = \frac{4\pi A_{\mathrm{eff}}}{\lambda^2}.
  \label{eq:gain-area}
\end{equation}
For a parabolic dish $A_{\mathrm{eff}} = \eta_A\,\pi D^2/4$; the
half-power beam width is
\begin{equation}
  \theta_{\mathrm{HPBW}} = k_b\,\frac{\lambda}{D},
\end{equation}
with $k_b = 1.02$ for uniform illumination and $k_b \approx 1.2$ for a
typical tapered feed (used in the simulations). For an antenna specified
only by its gain (e.\,g.\ a horn), the beam solid angle is
$\Omega_A = 4\pi/G$; with a Gaussian main beam,
$\Omega_A = 1.133\,\theta_{\mathrm{HPBW}}^2$, so
\begin{equation}
  \theta_{\mathrm{HPBW}} = \sqrt{\frac{4\pi}{1.133\,G}} .
\end{equation}

\paragraph*{Antenna temperature.} A source of flux density $S$ raises
the antenna temperature by
\begin{equation}
  T_A = \frac{A_{\mathrm{eff}}\,S}{2k},
  \label{eq:ta}
\end{equation}
where the factor $1/2$ accounts for a single polarisation receiving an
unpolarised source. For a source larger than the beam only the part
inside the beam contributes.

\section{Sensitivity}

The receiver adds noise of temperature $T_{\mathrm{sys}}$. Expressed as
the flux density that would produce the same antenna temperature,
\begin{equation}
  \mathrm{SEFD} = \frac{2\,k\,T_{\mathrm{sys}}}{A_{\mathrm{eff}}} .
\end{equation}
The radiometer equation gives the noise of a total-power measurement
with bandwidth $\Delta\nu$ and integration time $\tau$:
\begin{equation}
  \sigma_T = \frac{T_{\mathrm{sys}}}{\sqrt{\Delta\nu\,\tau}} .
  \label{eq:radiometer}
\end{equation}
A strong source raises the total noise: if $T_A$ is not small compared
with $T_{\mathrm{sys}}$, $T_{\mathrm{sys}}$ must be replaced by
$T_{\mathrm{sys}} + T_A$. For the Sun observed with a 3\,m dish
($T_A \approx 730\,\mathrm{K}$ against $T_{\mathrm{sys}} = 120\,\mathrm{K}$)
the receiver must also stay linear, which requires an attenuator in
front of the LNA.

\section{Spectral lines}

Interferosim supports the following lines (rest frequencies, see
\texttt{interferosim.constants}):
\begin{center}
\begin{tabular}{@{}lll@{}}
\toprule
Line & Frequency [MHz] & Remark \\
\midrule
HI 21\,cm & 1420.405\,751\,768 & hyperfine transition of neutral hydrogen \\
OH & 1612.231 & satellite line \\
OH & 1665.402 & main line \\
OH & 1667.359 & main line \\
OH & 1720.530 & satellite line \\
\bottomrule
\end{tabular}
\end{center}
